\title{Schwa to /\protect\MakeLowercase{a}/: Semi-longitudinal development of /\protect\MakeLowercase{a}/ in L1 English L2 Spanish learners}
\author{Robert Esposito, Joseph V. Casillas}
\organization{Department of Spanish and Portuguese, Rutgers University}
\email{rme70@scarletmail.rutgers.edu, joseph.casillas@rutgers.edu}


\maketitle

\begin{abstract}
Native (L1) English second language (L2) Spanish learners acquiring the Spanish vowel system must adjust entrenched vowel categories and morphophonological rules. English exhibits stress-conditioned allophonic variation nonexistent in Spanish, despite the triggering phonetic environments being present, which learners must come to disregard. This semi-longitudinal study analyzes Spanish /a/ in stressed and unstressed syllables produced in bisyllabic paroxytone nonce words by 10 L1 English L2 Spanish learners during a seven-week domestic immersion program. F1 and F2 were extracted from seven points of /a/ productions and averaged to obtain spectral centroids. Generalized additive mixed models revealed lower F1 and higher F2 in unstressed than stressed /a/. Over time, unstressed /a/ F1 rose to approach stressed /a/, while F2 increased in general, indicating phonological and phonetic development, respectively. These findings provide semi-longitudinal evidence of early L1 English L2 Spanish acquisition of phonological and phonetic properties of the Spanish vowel system.
\end{abstract}

\keywords{phonetics, phonology, second language acquisition, Spanish, English}

